\begin{savequote}[75mm]
What we have seen in financial markets should bring home to us all that the central organizing principle of this 21st century is interdependence.
\qauthor{Kevin Rudd}
\end{savequote}

\chapter{Introduction}

In the era of globalization, financial markets have become increasingly interconnected. Economies are linked by multinational companies and global supply chains. Economic conditions in one region of the world typically also influence others around the world. Financial shocks in one country can quickly impact markets in other countries, increasing the complexity of markets.

As global interconnectedness increases, understanding the mechanisms by which national markets are connected is an important task for investors, policymakers, and financial researchers. Investors wanting to diversity their portfolio may question to what degree markets can be independent of each other. Policymakers must understand the impact of international regulations on the markets. Financial researchers attempting to understand markets must now consider links with other international markets, increasing the complexity of the research.

\section{Background Information}

The 1980s and 1990s were fueled by market liberalization, in which many countries relaxed capital regulations and allowed money to flow more easily across borders.\cite{followingthemoney} This caused investors to start investing in foreign stocks and bonds, leading to a period of rapid growth in international equity and debt markets.\cite{followingthemoney} In the 1970s, the annual foreign direct investment by multinational corporations was about \$13 billion; by 2023, it had increased to \$1.37 trillion. The result of this era was the normalization of global portfolios and significantly more interconnected financial markets.

In the resulting markets, price movements, trends, and financial crises in one country's market can quickly advance to other markets. This is seen at a large scale in cases as early as the 1997 Asian Financial Crisis, during which a currency crisis in Thailand led to steep market declines in Malaysia, Indonesia, and South Korea and led to spillover effects in Russia, Latin America, and Eastern Europe.\cite{scenarios} The 2008 Financial Crisis began as the collapse of the U.S. housing bubble and evolved into the worst global recession in decades, leading to the major financial institutions across Europe requiring government bailouts within a year.\cite{scenarios} In 2015, a sharp selloff in the Chinese market resulted in the Dow Jones falling by over 1,000 points at the next trade open.\cite{scenarios} Most recently with the COVID-19 Pandemic in early 2020, uncertainty and fear by the pandemic sparked market crashes all around the world.\cite{scenarios}

More importantly, studies have shown that spillovers between the U.S. market and other regions increased after these global crises. The markets are so interconnected that they cannot live in isolation, and their interconnectedness will only keep growing.

The structure of this interdependence is an important field of study in international finance, typically focusing on links between two specific countries. Much of the literature has focused on measuring the integration of markets through correlations, co-movements, and global factors. We do not fully understand how information from one market can be used to statistically replicate the return structure of another. 

These strategies are especially important in the case of the U.S. and Chinese equity markets, which make up over 60\% of the global stock market capitalization. Historically, U.S. market signals influenced markets in Europe and Asia, with China remaining isolated and not impacted by foreign changes. This has changed dramatically over the past decade as China has now integrated itself in the global markets. China allowed international investors to trade certain Chinese A-shares for the first time in 2014.\cite{IMF} More than a decade later, U.S. investors hold increasing stakes in Chinese companies and Chinese firms are listed on American exchanges, demonstrating an interconnectedness that impacts both equity markets.

This thesis addresses the information gap by investigating the extent to which U.S. and Chinese macroeconomic factors and market-based signals can be used to replicate returns in the U.S. and Chinese equity markets. As a measure of economic co-integration, this thesis employs methods to estimate how much additional explanatory power is gained by adding one country's economic factors to a model that estimates the other country's market returns using it's own economic factors. If the additional set of factors add a lot of explanatory power, then there must be some level of market spillovers, with spillovers in both directions implying co-integration. This analysis is conducted using linear models with regularization, generalized additive models, and kernel regression. The investigations in this thesis will also include data from 13 themes and over 150 factors representing the U.S. and Chinese economies, with some measurements dating back to the 1920s. The data is also available at daily and monthly granularity, with 1,164 rows in the monthly dataset and nearly 3 million rows in the daily dataset.

\section{Literature Review}

\subsection{Foundational Asset Pricing Theories}

Introduced by Harry Markowitz in 1952, the Markowitz model was a revolutionary portfolio optimization model that maximized expected return at set levels of risk, typically measured by standard deviation.\cite{markowitz} The model showed that by diversifying a portfolio, investors can take advantage of a fronteir of optimal risk-return trade-offs. This model was so impactful that it started an entire new field: Modern Portfolio Theory.

The Capital Asset Pricing Model was introduced independently by Treynor, Sharpe, Lintner, and Mossin between 1961 and 1966.\cite{CAPM} It extended the pricing model by separating systematic risk from idiosyncratic risk. In 1964 and 1965, Sharpe and Lintner showed that the market portfolio, a portfolio in which each asset is weighted proportionally to its presence in the market, is mean-variance efficient under a few assumptions.\cite{CAPM} This implied that all investors should hold a combination of the market and a risk-free asset. 

However, the model did not hold empirically, and in 1976, Ross developed the Arbitrage Pricing Theory, which generalized this from a singular market risk factor to multiple risk factors.\cite{ross} The Arbitrage Pricing Theory does not specify the factors a priori and assumes that the relationship between expected returns and factor exposures is linear.

In the 1990s, Fama and French showed that a couple systematic factors can help explain movements in equity returns better than any single market beta.\cite{fama_french1992} The factors were initially size and value, though they added market in their three-factor model. Carhart introduced momentum as a fourth factor in 1997, and Fama and French proposed a five-factor model in 2015 with profitability and investment.\cite{carhart, fama_french2015} This has sparked an interest in using factors to understand equity market returns.

\subsection{Global Market Integration}

Whether or not a set of economic factors from one country can replicate the returns depends crucially on the degree to which the markets are globally integrated. If they are highly integrated, they can be impacted by the same global factors. Otherwise, the markets might be more impacted by local idiosyncratic factors.

In 1995, Bekaert and Harvey found time-varying world market integration among the 12 emerging markets they examined, many of which exhibited only partial integration.\cite{bekaert} They noted that integration is expected to increase with financial liberalization. During the late 20th century, developed markets in the West were fairly integrated and emerging markets in Asia and Latin America were more segmented.

Pukthuanthong and Roll showed in 2009 that correlation alone is not a valid measure of integration, arguing that markets that are perfectly integrated can have low correlations if they are impacted differently by global factors.\cite{pukthuanthong} They proposed using measure of the variance explained by models using the factors as predictors to predict the returns of the market—they specifically suggested $R^2$ in multi-factor models.\cite{pukthuanthong}

It's important to note that China had a largely isolationist policy until 2014, historically operating with tighter capital controls. As such, China was for a long time partially segmented from the world, and its markets were no different. Analyses done in the 2000s revealed little integration with the global economy.

\subsection{Cross-Market Return Replication as a Measure of Integration}

An emerging approach in the literature is to assess market integration by testing how well one market's return can be replicated using factors from another market. If two markets are integrated, then the economic factors that drive one market should explain some variability in the other. 

An attempt to explain Chinese A-share portfolio returns using U.S. risk factors for 1993-2006 data by Brooks et al. (2010) found that the Chinese market was almost completely segmented, with U.S. factors adding no explanatory power over the Chinese factors.\cite{brooks} As China's markets liberalized, replication studies found stronger links between the U.S. and Chinese markets. Goh et al. (2013) found no link between the countries before China entered the World Trade Organization in 2001 and a strong link after, characterized by a set of U.S. economic factors showing significant predictive power for Chinese stock returns.\cite{goh} This trend is increasing, with recent literature showing that the Chinese market becoming so integral to the global economy that its lagged returns can predict returns in numerous other markets worldwide. Chinese market integration is expected to continue to grow, extending trends from the past decade.

\subsection{Comparison with Traditional Integration Metrics}

Classical approaches for market integration metrics have relied on correlational studies bewteen the markets, cointegration testing for long-term relationships, and comparing factor exposures in global CAPM or mutli-factor models.

\subsubsection{Correlation-Based Measures}

Although Pukthuanthong and Roll showed that correlation is not a valid measure of integration, Billio et al. found that many advanced integration measures yielded the same patterns as rolling correlations.\cite{billio} Regardless of effectiveness, correlation-based measures only analyze co-movement strength and do not dive into the underlying drivers behind the movements of each market. Correlational studies also assume linear and uniform relationships between the markets. 

\subsubsection{Cointegration and Long-Run Parity Tests}

Another common approach is to test whether markets are cointegrated, typically assessed by the Johansen test. For two markets to be co-integrated, their divergences must revert to the mean. Though cointegration studies between China and the U.S. show cointegration only after China's market started liberalizing, cointegration is often a binary indicator with no value representing the strength of the relationship. 

\subsubsection{Shared Factor Exposure Models}

The most recent methodology is to replicate each market using a subset of global factors and compare the two sets of factors that explain the most variability in the returns. If the same global CAPM or multi-factor model can represent both markets without needing country-specific data, integration between the markets is implied. The Fama-French models are a popular attempt to do this with preset factors, often comprising of three or five factors. These are often limiting as they are a static set of factors that are limited in complexity. 

\subsection{Advantages of the Return Replication Approach}

Returns replication has several advantages over these methods. First, it analyzes the underlying effects of the integration and determines the relative strength of different factors. Second, the lack of a standardized set of factors allows for a more robust measurement that is more capable of capturing the true market relationship. The results naturally follow a practical economic meaning. Lastly, this replication approach is capable of taking advantage of modern econometric and machine learning techniques like linear models with regularization, generalized additive models, and kernel regressions.